\chapter{Interpolation and Polynomial Approximation}

\section{The Interpolation Problem}

Suppose that $n+1$ data points
\[
 (x_0,y_0),(x_1,y_1),\ldots,(x_n,y_n)
\]
are given, where the nodes $x_i$ are distinct.  Interpolation seeks a function $P$ satisfying
\[
 P(x_i)=y_i,
 \qquad i=0,1,\ldots,n.
\]
In polynomial interpolation, $P$ is chosen from
\[
 \mathcal P_n
 =\left\{a_nx^n+a_{n-1}x^{n-1}+\cdots+a_1x+a_0\right\}.
\]
Unlike a regression curve, an interpolating polynomial passes exactly through every data point.

\begin{theorem}[Weierstrass approximation theorem]
If $f$ is continuous on $[a,b]$, then for every $\varepsilon>0$ there exists a polynomial $P$ such that
\[
 |f(x)-P(x)|<\varepsilon
 \qquad\text{for all }x\in[a,b].
\]
\end{theorem}

The theorem guarantees that continuous functions can be approximated uniformly by polynomials, but it does not prescribe the degree or the best choice of nodes.

\begin{theorem}[Existence and uniqueness of the interpolating polynomial]
Given $n+1$ distinct nodes $x_0,\ldots,x_n$ and arbitrary values $y_0,\ldots,y_n$, there exists exactly one polynomial $P_n\in\mathcal P_n$ satisfying
\[
 P_n(x_i)=y_i,
 \qquad i=0,\ldots,n.
\]
\end{theorem}

\section{Cardinal and Lagrange Polynomials}

For each node $x_k$, define the cardinal polynomial
\[
 L_k(x)=\prod_{\substack{i=0\\i\ne k}}^n
 \frac{x-x_i}{x_k-x_i}.
\]
These polynomials satisfy
\[
 L_k(x_i)=
 \begin{cases}
 1,&i=k,\\
 0,&i\ne k.
 \end{cases}
\]
Therefore the Lagrange interpolating polynomial is
\[
 P_n(x)=\sum_{k=0}^n y_kL_k(x).
\]
At a node $x_i$, every term vanishes except $y_iL_i(x_i)=y_i$.

\begin{example}
For the data
\[
\begin{array}{c|ccc}
\toprule
x&1&3&4\\
\midrule
f(x)&2&5&8\\
\bottomrule
\end{array}
\]
find the Lagrange interpolating polynomial and estimate $f(2.5)$.
\end{example}

\begin{solution}[7.0cm]
The cardinal polynomials are
\[
 L_0(x)=\frac{(x-3)(x-4)}{(1-3)(1-4)}
       =\frac{(x-3)(x-4)}6,
\]
\[
 L_1(x)=\frac{(x-1)(x-4)}{(3-1)(3-4)}
       =-\frac{(x-1)(x-4)}2,
\]
\[
 L_2(x)=\frac{(x-1)(x-3)}{(4-1)(4-3)}
       =\frac{(x-1)(x-3)}3.
\]
Hence
\[
 P_2(x)=2L_0(x)+5L_1(x)+8L_2(x).
\]
After simplification,
\[
 P_2(x)=\frac12x^2-\frac12x+2.
\]
Therefore
\[
 f(2.5)\approx P_2(2.5)
 =\frac12(2.5)^2-\frac12(2.5)+2
 =3.8750.
\]
\end{solution}

\begin{remark}
The Lagrange form is symmetric in the nodes and conceptually direct.  However, if a new node is added, most of the basis polynomials must be rebuilt.  Newton's form is usually more convenient for incremental computation.
\end{remark}

\section{Divided Differences and Newton's Formula}

The Newton form of an interpolating polynomial is
\[
 P_n(x)=c_0+c_1(x-x_0)+c_2(x-x_0)(x-x_1)+\cdots
 +c_n\prod_{j=0}^{n-1}(x-x_j).
\]
The coefficients are divided differences.

\begin{definition}[Divided differences]
For distinct nodes, define
\[
 f[x_i]=f(x_i),
\]
\[
 f[x_i,x_{i+1}]
 =\frac{f[x_{i+1}]-f[x_i]}{x_{i+1}-x_i},
\]
and recursively
\[
 f[x_i,\ldots,x_{i+k}]
 =\frac{f[x_{i+1},\ldots,x_{i+k}]
       -f[x_i,\ldots,x_{i+k-1}]}
 {x_{i+k}-x_i}.
\]
\end{definition}

The Newton divided-difference polynomial is
\[
\begin{aligned}
P_n(x)={}&f[x_0]
+f[x_0,x_1](x-x_0)\\
&+f[x_0,x_1,x_2](x-x_0)(x-x_1)+\cdots\\
&+f[x_0,\ldots,x_n]
\prod_{j=0}^{n-1}(x-x_j).
\end{aligned}
\]

\begin{example}
For the points $(1,2)$, $(3,5)$ and $(4,8)$, compute
$f[1]$, $f[1,3]$ and $f[1,3,4]$, and construct the Newton interpolating polynomial.
\end{example}

\begin{solution}[5.5cm]
The first divided differences are
\[
 f[1]=2,
 \qquad
 f[1,3]=\frac{5-2}{3-1}=\frac32,
\]
\[
 f[3,4]=\frac{8-5}{4-3}=3.
\]
The second divided difference is
\[
 f[1,3,4]
 =\frac{f[3,4]-f[1,3]}{4-1}
 =\frac{3-\frac32}{3}
 =\frac12.
\]
Thus
\[
 P_2(x)=2+\frac32(x-1)+\frac12(x-1)(x-3).
\]
This simplifies to
\[
 P_2(x)=\frac12x^2-\frac12x+2,
\]
the same unique interpolating polynomial obtained in Lagrange form.
\end{solution}

\section{Forward and Backward Differences}

When the nodes are equally spaced,
\[
 x_i=x_0+ih,
\]
divided differences can be represented efficiently by finite-difference tables.

\begin{definition}[Forward difference]
The forward difference operator $\Delta$ is defined by
\[
 \Delta f(x)=f(x+h)-f(x),
\]
and recursively
\[
 \Delta^k f(x)=\Delta^{k-1}f(x+h)-\Delta^{k-1}f(x).
\]
\end{definition}

\begin{definition}[Backward difference]
The backward difference operator $\nabla$ is defined by
\[
 \nabla f(x)=f(x)-f(x-h),
\]
and recursively
\[
 \nabla^k f(x)=\nabla^{k-1}f(x)-\nabla^{k-1}f(x-h).
\]
\end{definition}

For equally spaced nodes,
\[
 f[x_0,x_1,\ldots,x_k]
 =\frac{\Delta^k f(x_0)}{k!h^k},
\]
and
\[
 f[x_n,x_{n-1},\ldots,x_{n-k}]
 =\frac{\nabla^k f(x_n)}{k!h^k}.
\]

\subsection{Newton Forward Difference Formula}

Let
\[
 s=\frac{x-x_0}{h}.
\]
Then
\[
 P_n(x)=f(x_0)
 +s\Delta f(x_0)
 +\frac{s(s-1)}{2!}\Delta^2f(x_0)
 +\cdots
 +\binom{s}{n}\Delta^nf(x_0),
\]
where
\[
 \binom{s}{k}
 =\frac{s(s-1)\cdots(s-k+1)}{k!}.
\]
This form is most convenient when $x$ is near the beginning of the table.

\subsection{Newton Backward Difference Formula}

Let
\[
 s=\frac{x-x_n}{h}.
\]
Then
\[
 P_n(x)=f(x_n)
 +s\nabla f(x_n)
 +\frac{s(s+1)}{2!}\nabla^2f(x_n)
 +\cdots
 +\frac{s(s+1)\cdots(s+n-1)}{n!}\nabla^nf(x_n).
\]
This form is most convenient when $x$ is near the end of the table.

\begin{example}
Consider the data
\[
 (1,0.76519),\quad
 (1.3,0.62008),\quad
 (1.6,0.45541),\quad
 (1.9,0.28182).
\]
\begin{enumerate}[label=(\alph*),leftmargin=2.7em]
  \item Use Newton's forward formula to estimate $f(1.1)$.
  \item Use Newton's backward formula to estimate $f(1.8)$.
\end{enumerate}
\end{example}

\begin{solution}[10.0cm]
Here $h=0.3$.  The forward-difference table begins
\[
\begin{array}{c|r|r|r|r}
\toprule
x&f&\Delta f&\Delta^2f&\Delta^3f\\
\midrule
1.0&0.76519&-0.14511&-0.01956&0.01064\\
1.3&0.62008&-0.16467&-0.00892&\\
1.6&0.45541&-0.17359&&\\
1.9&0.28182&&&\\
\bottomrule
\end{array}
\]
For $x=1.1$,
\[
 s=\frac{1.1-1}{0.3}=\frac13.
\]
Thus
\[
\begin{aligned}
P(1.1)={}&0.76519
+\frac13(-0.14511)
+\frac{(1/3)(-2/3)}{2}(-0.01956)\\
&+\frac{(1/3)(-2/3)(-5/3)}{6}(0.01064)
\approx0.71965.
\end{aligned}
\]
For the backward formula at $x=1.8$,
\[
 s=\frac{1.8-1.9}{0.3}=-\frac13.
\]
Using
\[
 \nabla f(1.9)=-0.17359,
 \quad
 \nabla^2f(1.9)=-0.00892,
 \quad
 \nabla^3f(1.9)=0.01064,
\]
we obtain
\[
\begin{aligned}
P(1.8)={}&0.28182
-\frac13(-0.17359)
+\frac{(-1/3)(2/3)}2(-0.00892)\\
&+\frac{(-1/3)(2/3)(5/3)}6(0.01064)
\approx0.340018.
\end{aligned}
\]
Therefore
\[
 f(1.1)\approx0.71965,
 \qquad
 f(1.8)\approx0.340018.
\]
\end{solution}

\begin{remark}
High-degree interpolation at equally spaced nodes can oscillate severely near the endpoints, a phenomenon known as Runge's phenomenon.  Increasing the number of nodes does not automatically improve the approximation everywhere.  Piecewise polynomials and carefully chosen nodes are often preferable.
\end{remark}
